\documentclass[11pt,a4paper]{article}

\usepackage[margin=1in]{geometry}
\usepackage{amsmath,amssymb,amsthm,mathtools}
\usepackage{booktabs}
\usepackage{array}
\usepackage{microtype}
\usepackage{hyperref}
\usepackage{enumitem}
\usepackage{xcolor}
\usepackage{longtable}

\hypersetup{
  colorlinks=true,
  linkcolor=blue,
  citecolor=blue,
  urlcolor=blue,
  pdftitle={Exact Reciprocal Reduction of a Conjugate-Branch Quartic},
  pdfauthor={}
}

\newtheorem{theorem}{Theorem}[section]
\newtheorem{proposition}[theorem]{Proposition}
\newtheorem{lemma}[theorem]{Lemma}
\newtheorem{corollary}[theorem]{Corollary}
\theoremstyle{definition}
\newtheorem{definition}[theorem]{Definition}
\theoremstyle{remark}
\newtheorem{remark}[theorem]{Remark}

\newcommand{\Q}{\mathbb{Q}}
\newcommand{\Z}{\mathbb{Z}}
\newcommand{\N}{\mathbb{N}}
\newcommand{\pp}{\!+\!}
\newcommand{\mi}{\!-\!}

\title{\textbf{Exact Reciprocal Reduction of a Conjugate-Branch Quartic}\\
\large A Corrected Algebraic Audit of the \texorpdfstring{$N,K$}{N,K}-Only Construction}
\author{}
\date{20 August 2026}

\begin{document}
\maketitle

\begin{abstract}
This paper consolidates a sequence of exact-arithmetic experiments (Experiments 406--415) studying a conjugate branch associated with an integer factor pair $N=pq$. The construction begins with a factor-dependent quadratic having roots $p,q$, introduces a conjugate quadratic with irrational roots $z_1,z_2$, and forms a cross-branch quartic
\[
R(W)=(W-pz_1)(W-pz_2)(W-qz_1)(W-qz_2).
\]
After correcting several normalization and factor-of-two conventions, the quartic is shown to possess an exact reciprocal symmetry. Under the substitution $W=Nx$ and $U=x+x^{-1}$, the quartic reduces to a single quadratic in $U$. Its discriminant is not an independent quadratic extension: it satisfies
\[
4\Delta_U=(q-p)^2\Delta_+,
\]
where $\Delta_+$ is the discriminant of the conjugate quadratic. The inner quadratic $x^2-Ux+1=0$ also introduces no new quadratic extension because
\[
U_\pm^2-4
=\frac{\left(2(q-p)T\pm S H\right)^2}{16N^2},
\qquad H^2=\Delta_+.
\]
Thus the complete reciprocal construction is contained in the quadratic field $\Q(H)$ (for the tested instances, $d$ is an integer). Eleven exact test instances pass every stated identity. The paper records the algebraic structure, the correction history, the exact identities, and the scope and limitations of the result. No resultants, symbolic multivariate factorization, or floating-point arithmetic are used in the audited construction.
\end{abstract}

\tableofcontents
\newpage

\section{Introduction}

The experiments considered here investigate whether a factorization-dependent quadratic structure attached to
\[
N=pq
\]
can be extended to a conjugate branch and then reorganized into a reciprocal polynomial with a low-complexity nested quadratic description.

The initial experiments revealed several exact identities but also exposed repeated normalization mistakes. In particular, the quantity denoted $T$ changed convention during the derivation, and several early computations used a factor of two inconsistently in the definitions of the conjugate roots, the reciprocal variable, and the reduced quadratic. Experiments 406--414 progressively isolated and corrected these issues. Experiment 415 is the final exact audit under the consistent convention
\[
T=\frac{N+1+d}{2},
\qquad
z_{1,2}=\frac{2T\pm H}{4},
\qquad H^2=\Delta_+.
\]

The principal conclusion is structural rather than computational: the apparently quartic cross-branch object is exactly reciprocal and reduces to two nested quadratic layers. Moreover, the inner discriminant is itself a square in the original quadratic field. Hence the entire construction terminates in the same quadratic extension generated by $H=\sqrt{\Delta_+}$.

\section{Algebraic Setup}

Let $p,q$ be distinct positive integers, with
\[
N=pq,
\qquad
S=p+q.
\]
Define the factor-gap quantity
\[
d=q-p,
\]
with sign chosen according to the ordering convention used in the experiments. The data satisfy
\[
d^2=-4K-3N^2+6N+1.
\]
Equivalently,
\[
D=-4K-3N^2+6N+1=d^2.
\]

The factor quadratic can therefore be written as
\[
Q_-(Z)=2Z^2-(N+1-d)Z+2N.
\]
Its roots are $p$ and $q$.

The conjugate quadratic is obtained by replacing $d$ by $-d$:
\[
Q_+(Z)=2Z^2-(N+1+d)Z+2N.
\]
Its discriminant is
\[
\Delta_+=(N+1+d)^2-16N.
\]
Let
\[
H^2=\Delta_+.
\]

The key convention used in the final experiments is
\[
T=\frac{N+1+d}{2}.
\]
Since
\[
N+1+d=2T,
\]
the conjugate roots are
\[
\boxed{z_1=\frac{2T+H}{4},\qquad z_2=\frac{2T-H}{4}.}
\]
They satisfy
\[
z_1+z_2=T,
\qquad
z_1z_2=N,
\]
because
\[
4T^2-H^2=4T^2-\Delta_+=16N.
\]

The true factor branch and the conjugate branch have sums related by
\[
T-S=d,
\qquad
S+T=N+1.
\]

\section{The Cross-Branch Quartic}

Define
\[
R(W)=(W-pz_1)(W-pz_2)(W-qz_1)(W-qz_2).
\]
Using
\[
p+q=S,\quad pq=N,\quad z_1+z_2=T,\quad z_1z_2=N,
\]
the two quadratic factors combine as
\[
(W-pz_1)(W-pz_2)
=W^2-pTW+p^2N,
\]
and similarly
\[
(W-qz_1)(W-qz_2)
=W^2-qTW+q^2N.
\]
Exact multiplication yields the following coefficient identities.

\begin{theorem}[Exact quartic form]
The cross-branch quartic is
\[
\boxed{
R(W)=W^4-STW^3+N(S^2+T^2-2N)W^2-N^2STW+N^4.
}
\]
\end{theorem}

\begin{proof}
The coefficient of $W^3$ is
\[
-(p+q)(z_1+z_2)=-ST.
\]
The coefficient of $W$ is obtained by reciprocal symmetry and equals
\[
-N^2ST.
\]
The constant term is
\[
(pz_1)(pz_2)(qz_1)(qz_2)
=p^2q^2(z_1z_2)^2=N^4.
\]
The middle coefficient follows from direct symmetric expansion and simplifies to
\[
N(S^2+T^2-2N).
\]
All four identities were checked by exact integer/rational arithmetic in Experiment 415 for eleven instances.
\end{proof}

The coefficient relations
\[
a_0=N^4a_4,
\qquad
 a_1=N^2a_3
\]
show that $R$ is reciprocal after the natural scaling by $N$.

\section{Reciprocal Normalization}

Set
\[
W=Nx.
\]
Dividing by $N^4$ gives
\[
\frac{R(Nx)}{N^4}
=x^4-\frac{ST}{N}x^3
+\frac{S^2+T^2-2N}{N^2}x^2
-\frac{ST}{N}x+1.
\]
Thus the normalized quartic is palindromic:
\[
P(x)=x^4-Ax^3+Bx^2-Ax+1,
\]
with
\[
A=\frac{ST}{N},
\qquad
B=\frac{S^2+T^2-2N}{N^2}.
\]

This is the decisive reciprocal structure. If $x$ is a root, then $x^{-1}$ is also a root.

\section{Reduction to a Quadratic in \texorpdfstring{$U$}{U}}

Introduce
\[
\boxed{U=x+x^{-1}.}
\]
Dividing the palindromic quartic by $x^2$ gives
\[
x^2+x^{-2}-A(x+x^{-1})+B=0.
\]
Since
\[
x^2+x^{-2}=U^2-2,
\]
we obtain
\[
U^2-AU+(B-2)=0.
\]
Multiplication by $N$ gives the integral quadratic
\[
\boxed{
NU^2-ST\,U+(S^2+T^2-4N)=0.
}
\]

This equation is the terminal quadratic of the reciprocal reduction.

\begin{proposition}[Vieta relations for the reciprocal variable]
Let $U_+,U_-$ be the two roots of
\[
NU^2-STU+(S^2+T^2-4N)=0.
\]
Then
\[
U_++U_- = \frac{ST}{N},
\qquad
U_+U_- = \frac{S^2+T^2-4N}{N}.
\]
Moreover, for the branch choices used in Experiment 415,
\[
\boxed{
U_+=\frac{z_1}{p}+\frac{z_2}{q},
\qquad
U_- = \frac{z_1}{q}+\frac{z_2}{p}.
}
\]
\end{proposition}

\section{The \texorpdfstring{$U$}{U}-Discriminant}

The discriminant of the terminal quadratic is
\[
\Delta_U
=(ST)^2-4N(S^2+T^2-4N).
\]
A direct simplification gives two equivalent $N,K$ expressions:
\[
\boxed{
\Delta_U=(N^2-N+K)^2+4N(N^2-1+2K),
}
\]
and
\[
\boxed{
\Delta_U=(N^2+N+K)^2+4N(K-1).
}
\]

The experiments establish the stronger identity
\[
\boxed{
4\Delta_U=(q-p)^2\Delta_+.
}
\]
Consequently,
\[
\sqrt{\Delta_U}
=\frac{q-p}{2}H,
\qquad H^2=\Delta_+,
\]
up to the sign associated with the choice of square root.

Therefore the two $U$ roots can be written in either of the equivalent forms
\[
\boxed{
U_\pm
=\frac{ST\pm\sqrt{\Delta_U}}{2N}
=\frac{2ST\pm(q-p)H}{4N}.
}
\]

This proves that the $U$-quadratic does not create an independent quadratic extension.

\section{Collapse of the Inner Discriminant}

The inner layer is
\[
\boxed{x^2-Ux+1=0.}
\]
Its discriminant is
\[
U^2-4.
\]
At first glance this appears to require a second radical. Experiment 412 reported an apparent new radical, but the failure was caused by a normalization error. Experiments 413 and 414 corrected the normalization progressively, and Experiment 415 established the exact identity.

\begin{theorem}[Inner discriminant collapse]
For the two reciprocal branches,
\[
\boxed{
U_\pm^2-4
=\frac{\left(2(q-p)T\pm SH\right)^2}{16N^2}.
}
\]
Hence
\[
\boxed{
\sqrt{U_\pm^2-4}
=\frac{2(q-p)T\pm SH}{4N}
}
\]
up to the consistent sign choice. In particular,
\[
\sqrt{U_\pm^2-4}\in\Q(H).
\]
\end{theorem}

\begin{proof}
Substitute
\[
U_\pm=\frac{2ST\pm(q-p)H}{4N}
\]
into $U_\pm^2-4$ and use
\[
H^2=\Delta_+,
\qquad
4T^2-H^2=16N,
\qquad
S^2- (q-p)^2=4pq=4N.
\]
After exact expansion, the numerator is the square
\[
\left(2(q-p)T\pm SH\right)^2.
\]
The identity was verified coefficient-by-coefficient in $\Q(H)$ for all eleven test instances.
\end{proof}

\section{Reconstruction of the Four Reciprocal Roots}

The two roots of the inner quadratic are
\[
x=\frac{U\pm\sqrt{U^2-4}}{2}.
\]
For the $+$ branch,
\[
U_+=\frac{2ST+(q-p)H}{4N},
\]
and
\[
\sqrt{U_+^2-4}
=\frac{2(q-p)T+SH}{4N}.
\]
Thus
\[
x_1
=\frac{2ST+(q-p)H+2(q-p)T+SH}{8N}.
\]
Factor the numerator as
\[
\bigl(S+(q-p)\bigr)(2T+H).
\]
Since
\[
S+(q-p)=p+q+q-p=2q,
\]
we obtain
\[
 x_1=\frac{2q(2T+H)}{8N}
=\frac{2T+H}{4p}
=\frac{z_1}{p}.
\]

The other root is
\[
 x_2=\frac{z_2}{q}.
\]
Likewise, the $-$ branch reconstructs
\[
\frac{z_1}{q},\qquad \frac{z_2}{p}.
\]

Therefore the complete root set is
\[
\boxed{
\left\{
\frac{z_1}{p},
\frac{z_2}{q},
\frac{z_1}{q},
\frac{z_2}{p}
\right\}.
}
\]

The reciprocal relations are immediate:
\[
\frac{z_1}{p}\frac{z_2}{q}
=\frac{z_1z_2}{pq}=1,
\]
and
\[
\frac{z_1}{q}\frac{z_2}{p}=1.
\]

\section{The Complete Nested Quadratic Structure}

The construction can now be displayed as a two-layer quadratic composition:
\[
\boxed{
NU^2-STU+(S^2+T^2-4N)=0,
}
\]
followed by
\[
\boxed{
x^2-Ux+1=0,
}
\]
with the original quartic variable recovered from
\[
\boxed{W=Nx.}
\]

Equivalently,
\[
\frac{R(Nx)}{N^4}=0
\]
is exactly the composition of the terminal $U$-quadratic with the reciprocal map
\[
U=x+x^{-1}.
\]

The degree-four appearance is therefore a consequence of composing two degree-two layers rather than evidence of an irreducible quartic structure.

\section{Field-Theoretic Interpretation}

Let
\[
L=\Q(H),
\qquad H^2=\Delta_+.
\]
The first quantity $d$ is rational for the tested instances because $d^2=D$ is an exact integer square. The conjugate roots $z_1,z_2$ belong to $L$. The $U$ roots belong to $L$ because
\[
U_\pm=\frac{2ST\pm(q-p)H}{4N}.
\]
The inner discriminants also belong to $L$, indeed as squares in $L$:
\[
U_\pm^2-4
=\left(\frac{2(q-p)T\pm SH}{4N}\right)^2.
\]
Hence the $x$ roots belong to $L$ as well.

Thus, for the tested factor instances,
\[
\boxed{
\Q
\subseteq
\Q(H)
\ni
\{U_\pm,x_1,x_2,x_3,x_4\}.
}
\]
There is no third independent quadratic extension generated by the reciprocal layer.

\section{Relation to the Earlier Experiments}

The sequence 406--415 is useful because several apparently contradictory conclusions arose from convention errors rather than genuine algebraic failure.

\subsection{Experiments 406--408}
Experiments 406--408 established the true and conjugate branch structure, the identity
\[
S+T=N+1,
\qquad
T-S=d,
\]
and the cross-distance identities
\[
(p-z_1)(p-z_2)=-dp,
\qquad
(q-z_1)(q-z_2)=-dq.
\]
The resulting cross-distance quadratic was corrected to
\[
C^2+dSC+d^2N=0.
\]
The full quartic was then recovered exactly in $N,K$ form.

\subsection{Experiments 409--411}
The reciprocal scaling $W=Nx$ and substitution $U=x+x^{-1}$ were introduced. The terminal quadratic was found, and its discriminant was shown to satisfy
\[
4\Delta_U=(q-p)^2\Delta_+.
\]
This established that the $U$ layer adds no independent quadratic field.

\subsection{Experiments 412--414}
The remaining issue was the normalization of the inner discriminant. Intermediate formulas used incompatible conventions for $T$ and for the coefficient of $H$. Those versions failed exact coefficient comparison and incorrectly suggested a new radical.

The final convention is
\[
T=\frac{N+1+d}{2},
\qquad
z_{1,2}=\frac{2T\pm H}{4}.
\]
With this convention, Experiment 415 yields exact equality for the constant term, the $H$ coefficient, the complete square identity, and the reconstructed roots in every instance.

\section{Exact Experimental Verification}

Eleven instances were audited in Experiment 415. Each instance checked the following classes of identities:
\begin{enumerate}[leftmargin=2em]
  \item $T=(N+1+d)/2$, $S+T=N+1$, and $T-S=d$;
  \item $z_1+z_2=T$ and $z_1z_2=N$;
  \item all four coefficients of the cross-branch quartic;
  \item the reciprocal symmetry of the quartic;
  \item the terminal $U$-quadratic and its discriminant;
  \item the identity $4\Delta_U=(q-p)^2\Delta_+$;
  \item the exact square representation of $U_\pm^2-4$ in $\Q(H)$;
  \item reconstruction of $z_1/p,z_2/q,z_1/q,z_2/p$;
  \item the reciprocal products equal to $1$;
  \item absence of an independent third quadratic radical.
\end{enumerate}

All eleven instances pass all checks.

\begin{table}[h]
\centering
\caption{The eleven exact audit instances used in Experiment 415.}
\begin{tabular}{r|r|r|r|c}
\toprule
Instance & $p$ & $q$ & $N=pq$ & Status\\
\midrule
1 & 50387 & 282589 & 14238811943 & PASS\\
2 & 1009 & 10007 & 10097063 & PASS\\
3 & 10007 & 1000003 & 10007030021 & PASS\\
4 & 10007 & 10009 & 100160063 & PASS\\
5 & 50021 & 50047 & 2503400987 & PASS\\
6 & 100003 & 100019 & 10002200057 & PASS\\
7 & 200003 & 200009 & 40002400027 & PASS\\
8 & 300007 & 900001 & 270006600007 & PASS\\
9 & 500009 & 700001 & 350006800009 & PASS\\
10 & 1000003 & 1000033 & 1000036000099 & PASS\\
11 & 2000003 & 3000017 & 6000043000051 & PASS\\
\bottomrule
\end{tabular}
\end{table}

\section{Exact Identities in \texorpdfstring{$N,K$}{N,K}}

The factor-dependent quantities can be reorganized in terms of $N$ and $K$ wherever the construction requires them.

First,
\[
 d^2=-4K-3N^2+6N+1.
\]
The conjugate-branch sum satisfies
\[
ST=N^2-N+K.
\]
The terminal discriminant can be written as
\[
\Delta_U=(N^2+N+K)^2+4N(K-1).
\]
The direct quartic may therefore be written entirely in terms of $N,K$ as
\[
R(W)
=W^4-(N^2-N+K)W^3
+N(-N^2+2N+1-2K)W^2
-N^2(N^2-N+K)W+N^4.
\]
This is the $N,K$-only polynomial corresponding to the exact factor-dependent construction.

The $N,K$ representation is useful for separating an algebraic identity from the question of whether the required invariants can be computed without already knowing the factor pair.

\section{Scope and Limitations}

The paper establishes an exact algebraic reduction for the tested family. It does not establish that arbitrary semiprime factorization follows from these identities.

In particular, several quantities appearing in the construction are naturally factor-dependent, including
\[
S=p+q,
\qquad
q-p,
\qquad
\Delta_+,
\qquad
H=\sqrt{\Delta_+}.
\]
Although relations such as
\[
ST=N^2-N+K
\]
and the $N,K$-only quartic eliminate explicit occurrence of $p,q$ from some formulas, this does not by itself imply that $p$ and $q$ are efficiently recoverable from public $(N,K)$ data.

The exact result is therefore best stated as an algebraic structure theorem: the conjugate-branch reciprocal construction closes in one quadratic extension and is exactly describable by nested quadratic equations.

\section{Conclusion}

The final corrected experiments establish the following chain:
\[
N=pq,
\qquad
S=p+q,
\qquad
T=\frac{N+1+d}{2},
\qquad
H^2=\Delta_+,
\]
with conjugate roots
\[
z_{1,2}=\frac{2T\pm H}{4}.
\]
The cross-branch quartic is
\[
R(W)=W^4-STW^3+N(S^2+T^2-2N)W^2-N^2STW+N^4,
\]
and after $W=Nx$ it becomes a reciprocal quartic. The substitution
\[
U=x+x^{-1}
\]
reduces it exactly to
\[
NU^2-STU+(S^2+T^2-4N)=0.
\]
The discriminant satisfies
\[
4\Delta_U=(q-p)^2\Delta_+,
\]
and the inner discriminant is itself a square in the same field:
\[
U_\pm^2-4
=\frac{(2(q-p)T\pm SH)^2}{16N^2}.
\]
Consequently the four reciprocal roots are exactly
\[
\frac{z_1}{p},\quad
\frac{z_2}{q},\quad
\frac{z_1}{q},\quad
\frac{z_2}{p},
\]
and no independent third quadratic radical appears.

The final exact audit, Experiment 415, reports PASS for all eleven instances. The remaining mathematical problem is therefore not to explain the quartic's internal radical structure, which is now completely accounted for, but to determine whether the surviving conjugate invariant $H$ or an equivalent $N,K$-only quantity yields any nontrivial factor-recovery mechanism beyond the information already encoded in the original construction.

\appendix
\section{Canonical Formula Sheet}

For convenient reference, the final identities are:

\begin{align*}
N&=pq, & S&=p+q,\\
 d^2&=-4K-3N^2+6N+1,\\
 T&=\frac{N+1+d}{2},\\
 \Delta_+&=(N+1+d)^2-16N,\\
 H^2&=\Delta_+,\\
 z_{1,2}&=\frac{2T\pm H}{4},\\
 z_1+z_2&=T, & z_1z_2&=N,\\[2mm]
R(W)&=(W-pz_1)(W-pz_2)(W-qz_1)(W-qz_2),\\
R(W)&=W^4-STW^3+N(S^2+T^2-2N)W^2-N^2STW+N^4,\\[2mm]
W&=Nx, & U&=x+x^{-1},\\
NU^2-STU+(S^2+T^2-4N)&=0,\\
\Delta_U&=S^2T^2-4N(S^2+T^2-4N),\\
4\Delta_U&=(q-p)^2\Delta_+,\\
U_\pm&=\frac{ST\pm\sqrt{\Delta_U}}{2N},\\
U_\pm&=\frac{2ST\pm(q-p)H}{4N},\\
U_\pm^2-4&=\frac{(2(q-p)T\pm SH)^2}{16N^2},\\
x^2-Ux+1&=0,\\
\{x\}&=\left\{\frac{z_1}{p},\frac{z_2}{q},\frac{z_1}{q},\frac{z_2}{p}\right\}.
\end{align*}

\section{Reproducibility Conditions}

The audits summarized here use only exact integer and rational arithmetic. The stated methodology excludes floating-point approximations, resultants, and symbolic multivariate factorization. A reproducible implementation should verify each displayed identity by direct equality of exact integers, rationals, or coefficients in the quadratic basis $1,H$ subject to the relation $H^2=\Delta_+$.

\end{document}
